\documentclass[11pt]{article}

\usepackage[margin=1in]{geometry}
\usepackage[T1]{fontenc}
\usepackage[utf8]{inputenc}
\usepackage{textcomp}
\usepackage{amsmath}
\usepackage{amssymb}
\usepackage{booktabs}
\usepackage{array}
\usepackage{calc}
\usepackage{longtable}
\usepackage{xcolor}
\usepackage{algorithm}
\usepackage[noend]{algpseudocode}
\usepackage[hidelinks]{hyperref}
% Generated by scripts/paper/build_settings.jl; do not edit.
\providecommand{\pvDefine}[2]{\expandafter\def\csname pv@#1\endcsname{#2}}
\providecommand{\pv}[1]{\csname pv@#1\endcsname}
\pvDefine{settingAdamBetaOne}{0.9}
\pvDefine{settingAdamBetaTwo}{0.999}
\pvDefine{settingAdamEpsilon}{10^{-8}}
\pvDefine{settingAgentWidth}{16}
\pvDefine{settingBrokerDimension}{44}
\pvDefine{settingBrokerSeedCap}{100}
\pvDefine{settingBrokerWidth}{64}
\pvDefine{settingCalibrationMarkets}{50}
\pvDefine{settingCalibrationN}{1{,}000}
\pvDefine{settingCalibrationPairsGrouped}{10{,}000}
\pvDefine{settingDgpDimensionsWord}{eight}
\pvDefine{settingDgpProjectionWord}{three}
\pvDefine{settingE_init_agent}{100}
\pvDefine{settingE_init_broker}{100}
\pvDefine{settingEarlyWindow}{51--70}
\pvDefine{settingFrequencies}{1,2,3,4,5}
\pvDefine{settingFrequencyMax}{5}
\pvDefine{settingFrequencyMin}{1}
\pvDefine{settingHoldoutPartners}{40}
\pvDefine{settingHoldoutPrincipals}{100}
\pvDefine{settingIntervalPercent}{95}
\pvDefine{settingK}{5}
\pvDefine{settingLateWindow}{481--500}
\pvDefine{settingMinCandidatePairsWord}{five}
\pvDefine{settingN}{1000}
\pvDefine{settingNGrouped}{1{,}000}
\pvDefine{settingNestedBaseline}{20}
\pvDefine{settingNestedGeneral}{10}
\pvDefine{settingQ}{2.84}
\pvDefine{settingRosterSize}{200}
\pvDefine{settingSuppIntervalPercent}{95}
\pvDefine{settingT}{500}
\pvDefine{settingVarianceDivisor}{6}
\pvDefine{settingd}{8}
\pvDefine{settingdWord}{eight}
\pvDefine{settingdelta}{0.5}
\pvDefine{settingeta}{0.02}
\pvDefine{settingeta_lr_agent}{0.003}
\pvDefine{settingeta_lr_broker}{0.03}
\pvDefine{settingk}{6}
\pvDefine{settingn_strangers}{10}
\pvDefine{settingnetwork_measure_interval}{20}
\pvDefine{settingomega}{0.2}
\pvDefine{settingp_demand}{0.5}
\pvDefine{settingp_demandFixed}{0.50}
\pvDefine{settingp_rewire}{0.1}
\pvDefine{settingreservation_frac}{0.6}
\pvDefine{settingreservation_fracFixed}{0.60}
\pvDefine{settingrho}{0.5}
\pvDefine{settingrhoFixed}{0.50}
\pvDefine{settingridge_lambda_agent}{0.001}
\pvDefine{settingridge_lambda_broker}{0.001}
\pvDefine{settingroster_churn}{0.02}
\pvDefine{settingroster_frac}{0.2}
\pvDefine{settingroster_fracFixed}{0.20}
\pvDefine{settings}{8}
\pvDefine{settingsearch_cost_rate}{0.05}
\pvDefine{settingsigma_eps}{0.28}
\pvDefine{settingsigma_x}{0.5}
\pvDefine{settingtrain_max_obs}{2000}
\pvDefine{settingtrain_steps_agent}{50}
\pvDefine{settingtrain_steps_broker}{50}
\pvDefine{settingtrain_window_periods}{40}


\setlength{\emergencystretch}{3em}
\setlength{\parskip}{0.35em plus 0.1em minus 0.05em}
\setcounter{secnumdepth}{0}
\providecommand{\tightlist}{%
  \setlength{\itemsep}{0pt}\setlength{\parskip}{0pt}}
\ifdefined\DeclareUnicodeCharacter
  \DeclareUnicodeCharacter{00A7}{\S}
\fi
\newcommand{\R}{\mathbb{R}}

\title{Appendix B\\Model Specifications for an ABM of Brokerage in Matching Markets}
\author{}
\date{}

\begin{document}
\maketitle

This document specifies an agent-based model of brokered matching.

The model is a discrete-time simulation of a matching market with \emph{principals} and a \emph{broker}. Principals seeking pairwise matches either search on their own or outsource to the broker. All economic quantities (match output, fees, surplus) are in the same payoff units. All principals use heuristic decision rules. No principal holds beliefs about other principals' strategies.

The model specification covers principals (§1), the matching problem (§2), how principals learn to predict match value (§3), match economics (§4), network structure and principal turnover (§5), search and broker access (§6), the outsourcing decision (§7), and period measures (§8). Baseline parameter values and reproducibility are documented in Part II (§9).

\section{Part I. Model Specification}\label{part-i.-model-specification}

\subsection{1. Principals and Time}\label{agents-and-time}

The model runs for \(T\) discrete periods (default \pv{settingT}). It has \(N\) principals (default \pv{settingN}), a single broker, and an undirected network \(G\) over the principals and broker.

Each principal \(i\) is characterized by:

\begin{itemize}
\tightlist
\item \textbf{Type} \(\mathbf{x}_i \in \mathbb{R}^d\): a fixed vector of observable characteristics, drawn as specified below and used by the matching function (§2).
\item \textbf{Current-period matches} \(M_i^t\): active bilateral relationships involving \(i\) in period \(t\). Each counterparty can appear at most once in \(M_i^t\); incoming accepted offers are not capped.
\item \textbf{Active demand bound} \(K\): the maximum number of active demands a principal can initiate in a period (default \(K = \pv{settingK}\); §6).
\item \textbf{Experience history} \(\mathcal{H}_{i}^t = \{(\mathbf{x}_j, q_{ij})\}\): (other party's type, realized match output) pairs from all matches involving \(i\) (§3a).
\item \textbf{Satisfaction indices} \(s_{i,c}^t\): one scalar per search channel \(c \in \{\text{self}, \text{broker}\}\), used in the outsourcing decision (§7).
\end{itemize}

\subsubsection{Principal Types}\label{agent-types}

Type vectors are multidimensional but concentrated near a smooth one-dimensional curve on the unit sphere in \(\mathbb{R}^d\). Let \(u \in [0,1]\) index position along the curve:

\[\boldsymbol{\gamma}(u) = \frac{\tilde{\boldsymbol{\gamma}}(u)}{\|\tilde{\boldsymbol{\gamma}}(u)\|}, \qquad \tilde{\gamma}_\ell(u) = \begin{cases} \sin(2\pi \nu_\ell u + \psi_\ell) & \ell = 1, \ldots, d_\gamma \\ 0 & \ell = d_\gamma+1, \ldots, d. \end{cases}\]

The frequencies \(\nu_\ell \sim \operatorname{Unif}\{\pv{settingFrequencies}\}\) and phases \(\psi_\ell \sim \operatorname{Unif}[0,2\pi)\) are drawn once per simulation. The parameter \(d_\gamma\leq d\) is the number of coordinates in which the curve varies.

Each principal is drawn at a random position \(u_i \sim \mathrm{Unif}[0,1]\) on the curve, then perturbed:

\[\mathbf{x}_i = \frac{\boldsymbol{\gamma}(u_i) + \boldsymbol{\epsilon}_i}{\|\boldsymbol{\gamma}(u_i) + \boldsymbol{\epsilon}_i\|}, \qquad \boldsymbol{\epsilon}_i \sim \mathcal{N}\!\left(\mathbf{0}, \frac{\sigma_x^2}{d} \mathbf{I}_d\right).\]

Noise is applied in all \(d\) dimensions, then the vector is re-normalized. The scale \(\sigma_x / \sqrt{d}\) keeps the expected displacement from the curve approximately \(\sigma_x\) for any \(d\).

Thus the curve remains one-dimensional for every \(d_\gamma\). Increasing \(d_\gamma\) changes how many coordinates contain systematic curve variation; the remaining coordinates contain only perturbation noise.

\subsubsection{Broker}\label{broker}

The broker is a permanent node in \(G\) and has no type vector. Its state is:

\begin{itemize}
\tightlist
\item \textbf{Experience history} \(\mathcal{H}_b^t = \{(\mathbf{x}_i, \mathbf{x}_j, q_{ij})\}\): (party 1 type, party 2 type, realized match output) triples from mediated matches (§3c). The party labels have no directional meaning because the broker's pair features are symmetric.
\item \textbf{Roster} \(\text{Roster}^t\): the broker's standing access base (§6b).
\item \textbf{Current clients} \(D^t\): principals who outsource to the broker in period \(t\) (§6b).
\item \textbf{Reputation} \(\text{rep}_b^t\): the average broker satisfaction of current clients (§7).
\end{itemize}

\subsection{2. The Matching Problem}\label{the-matching-problem}

Each potential pair's underlying match value combines two components: the average of the principals' individual-quality scores and complementarity between their types. Each principal's individual-quality score is the same across partners, whereas complementarity depends on the pairing. The parameter \(\rho\) controls the relative contributions of these components. The parameter \(\delta\) makes complementarity harder to predict by increasing the difference between otherwise similar pairs on opposite sides of a boundary in type-pair space. A pair's underlying match value does not depend on which principal initiates the match. Both principals observe the same realized match output.

\subsubsection{2a. Match Output}\label{a.-match-output}

Let \(q_{ij}\) be the realized per-period output of a match between principals \(i\) and \(j\), measured in payoff units:

\[q_{ij} = Q + f(\mathbf{x}_i, \mathbf{x}_j) + \varepsilon_{ij}, \qquad
\varepsilon_{ij} \sim \mathcal{N}(0, \sigma_\varepsilon^2).\]

A pair's underlying match value is \(Q+f(\mathbf{x}_i,\mathbf{x}_j)\), determined by the two types and the matching-function objects fixed at initialization. Because the shock \(\varepsilon_{ij}\) is independent and has mean zero, underlying match value also equals \(E[q_{ij}\mid\mathbf{x}_i,\mathbf{x}_j]\), holding the matching environment fixed. The shock creates variation in realized match output that neither the principals nor the broker can predict.

The function \(f: \mathbb{R}^d \times \mathbb{R}^d \to \mathbb{R}\) is unknown to all principals:

\[
f(\mathbf{x}_i,\mathbf{x}_j)
=\frac{1}{s_{\mathrm{mix}}(\rho,\delta)}
\left[\sqrt{\rho}\,\widetilde U_{ij}
+\sqrt{1-\rho}\,\widetilde W_{ij}(\delta)\right].
\]

The tildes indicate that the individual-quality component \(U_{ij}\) (§2b) and the complementarity component \(W_{ij}(\delta)\) (§2c) are centered and scaled using fixed moments. These fixed moments are calculated across all unordered pairs of distinct principals within each of \pv{settingCalibrationMarkets} independently generated calibration markets of \pv{settingCalibrationN} principals, then averaged with equal weight across markets. The resulting coefficients depend on \(\rho\) and \(\delta\), are shared across simulations, and remain fixed through turnover.

Let \(\kappa(\delta)\) denote the average within-market covariance between the standardized individual-quality and complementarity components in the calibration markets. The common adjustment

\[s_{\mathrm{mix}}^2(\rho,\delta)=1+2\sqrt{\rho(1-\rho)}\,\kappa(\delta)\]

accounts for this covariance. Each standardized component and the combined systematic component \(f\) have unit average within-market variance in the calibration markets for every \(\rho\) and \(\delta\). Realized within-market variance of \(f\) can differ from one and change with turnover. The calibration normalization defines the payoff unit. In these units, \(Q=\pv{settingQ}\) and \(\sigma_\varepsilon=\pv{settingsigma_eps}\).

The standardized components have equal average within-market variances across the calibration markets. On this calibration scale, after weighting, the individual-quality component accounts for a share \(\rho\) of their summed marginal variances, and complementarity accounts for \(1-\rho\). At \(\rho=0\), match value depends only on complementarity; at \(\rho=1\), it depends only on individual quality. At \(\rho=1\), \(\delta\) has no effect.

\subsubsection{2b. Individual Quality}\label{b.-agent-individual-quality}

The vector \(\mathbf{c}\) defines a common ideal type and is fixed within each simulation. Principal \(i\)'s individual-quality score, \(\mathbf{x}_i^\top\mathbf{c}\), measures its alignment with that ideal. A pair's individual-quality component is the average of the two principals' scores:

\[
U_{ij}=\frac{1}{2}(\mathbf{x}_i^\top\mathbf{c}+\mathbf{x}_j^\top\mathbf{c}),
\qquad
\widetilde U_{ij}=\frac{U_{ij}-\mu_U}{\sigma_U}.
\]

The ideal type \(\mathbf{c}\in\mathbb{R}^d\) is drawn at initialization as a perturbation of a random point on the principal type curve:

\[u_c\sim\operatorname{Unif}[0,1],\qquad
\boldsymbol{\epsilon}_c\sim\mathcal N\!\left(\mathbf 0,\frac{\sigma_x^2}{d}\mathbf I_d\right),\qquad
\mathbf c=\frac{\boldsymbol{\gamma}(u_c)+\boldsymbol{\epsilon}_c}
{\left\|\boldsymbol{\gamma}(u_c)+\boldsymbol{\epsilon}_c\right\|_2}.\]

The normalization places \(\mathbf c\) on the same unit sphere as principal types. The perturbation therefore changes the ideal direction without changing its magnitude.

\subsubsection{2c. Match-Specific Complementarity}\label{c.-match-specific-complementarity}

Define base complementarity, signed complementarity, and difficulty-adjusted complementarity as

\[
V_{ij}=\mathbf{x}_i^\top\mathbf A\mathbf{x}_j,
\qquad
H_{ij}=\operatorname{sign}(\mathbf{x}_i^\top\mathbf B\mathbf{x}_j)V_{ij},
\qquad
W_{ij}(\delta)=V_{ij}+\delta H_{ij}.
\]

The normalized complementarity used in the matching function is

\[
\widetilde W_{ij}(\delta)
=\frac{W_{ij}(\delta)-\mu_W(\delta)}{\sigma_W(\delta)}.
\]

\(\mathbf{A}\) is a symmetric positive definite matrix that gives a smooth bilinear complementarity geometry. Because \(V_{ij}\) depends jointly on both types, it cannot in general be represented as the sum of separate principal scores. A second symmetric matrix, \(\mathbf{B}\), assigns pairs to two regions of type-pair space. Both matrices are fixed for the duration of the simulation.

\textbf{Base complementarity.} Draw \(\mathbf{M}_A\in\mathbb{R}^{d\times d}\) with iid \(\mathcal{N}(0,1)\) entries and set

\[
\mathbf{A}=\mathbf{M}_A^\top\mathbf{M}_A\cdot\frac{d}{\operatorname{tr}(\mathbf{M}_A^\top\mathbf{M}_A)}.
\]

Thus \(\mathbf{A}\) is symmetric positive definite and \(\operatorname{tr}(\mathbf{A})=d\).

The trace normalization prevents arbitrary scale differences in \(\mathbf{A}\) across realizations. The fixed pair-level normalization described in §2a then places complementarity on the common scale used in the matching function.

\textbf{Difficulty parameter.} Equivalently, write \(W_{ij}(\delta)=g(\mathbf{x}_i,\mathbf{x}_j)V_{ij}\), where a second symmetric matrix \(\mathbf{B}\in\mathbb{R}^{d\times d}\) defines

\[g(\mathbf{x}_i, \mathbf{x}_j) = 1 + \delta \cdot \text{sign}(\mathbf{x}_i^\top \mathbf{B} \mathbf{x}_j)\]

where \(\delta\in[0,1]\) (default \pv{settingdelta}).

The multiplier is \(1+\delta\) when \(\mathbf{x}_i^\top\mathbf{B}\mathbf{x}_j>0\), \(1-\delta\) when it is negative, and 1 when it is zero. At \(\delta=0\), the boundary between the two regions has no effect. For \(\delta>0\), pairs on opposite sides of the boundary can have sharply different underlying match values even when their types are otherwise similar. Increasing \(\delta\) increases this difference, making complementarity harder to learn. The scale \(\sigma_W(\delta)\) holds the average within-market variance of standardized complementarity at one in the calibration markets as \(\delta\) varies.

\textbf{Boundary construction.} At initialization, \(\mathbf{B}\) is made orthogonal to \(\mathbf{A}\) under the weighted inner product defined by the initial principals\textquotesingle{} empirical second-moment matrix. The resulting boundary remains fixed through turnover. Let

\[\boldsymbol{\Sigma}_x = \frac{1}{N} \sum_{i=1}^N \mathbf{x}_i \mathbf{x}_i^\top\]

be the uncentered second-moment matrix of the initial type vectors.

Draw \(\mathbf{Z}\in\mathbb{R}^{d\times d}\) with iid \(\mathcal{N}(0,1)\) entries and form a symmetric zero-trace matrix:

\[
\widetilde{\mathbf B}=\frac{\mathbf Z+\mathbf Z^\top}{2},
\qquad
\mathbf{B}_0 = \widetilde{\mathbf B} - \frac{\operatorname{tr}(\widetilde{\mathbf B})}{d}\mathbf{I}_d.
\]

Define the weighted inner product

\[\langle \mathbf{M}, \mathbf{N} \rangle_{\boldsymbol{\Sigma}_x} = \operatorname{tr}(\boldsymbol{\Sigma}_x \mathbf{M} \boldsymbol{\Sigma}_x \mathbf{N}),\]

then remove the projection of \(\mathbf{B}_0\) onto \(\mathbf{A}\) and normalize:

\[\mathbf{B}_{\text{raw}} = \mathbf{B}_0 - \frac{\operatorname{tr}(\boldsymbol{\Sigma}_x \mathbf{B}_0 \boldsymbol{\Sigma}_x \mathbf{A})}{\operatorname{tr}(\boldsymbol{\Sigma}_x \mathbf{A} \boldsymbol{\Sigma}_x \mathbf{A})} \mathbf{A}, \qquad \mathbf{B} = \frac{\mathbf{B}_{\text{raw}}}{\lVert \mathbf{B}_{\text{raw}} \rVert_F}.\]

This gives \(\operatorname{tr}(\boldsymbol{\Sigma}_x\mathbf{B}\boldsymbol{\Sigma}_x\mathbf{A})=0\) and \(\|\mathbf{B}\|_F=1\). Weighting by \(\boldsymbol{\Sigma}_x\) removes overlap between the two bilinear scores under the initial empirical type distribution, rather than in unused regions of \(\mathbb{R}^d\).

Symmetry of \(\mathbf{A}\) and \(\mathbf{B}\) makes \(f(\mathbf{x}_i,\mathbf{x}_j)=f(\mathbf{x}_j,\mathbf{x}_i)\).

\subsection{3. Learning}\label{learning}

Before searching for matches, principals and the broker may update their prediction models on their accumulated histories. Updated models are used to score candidates (§6). The broker is eligible to refit every period. Principal \(i\) is eligible when \(i\bmod 2=t\bmod 2\), so each principal is eligible every other period. An eligible model refits only when its history is nonempty and it has received new observations since its previous fitting call. A parameter switch selects the prediction model. Neural network (NN) learning is the default. Ridge regression is the alternate experiment.

\subsubsection{3a. Prediction Models and Fitting}\label{a.-architecture-and-fitting}

In the default model, both principals and the broker use a fully connected network with one hidden layer, ReLU activations, and a single linear output:

\[\hat{q}(\mathbf{z}) = \mathbf{w}_2^\top \text{ReLU}(\mathbf{W}_1 \mathbf{z} + \mathbf{b}_1) + b_2\]

where \(\hat q\) is predicted match value, \(\mathbf{z}\) is the input feature vector and the remaining symbols are network weights and biases. Principals receive raw partner type vectors; the broker receives pairs of them.

\textbf{Training rule.} Network weights minimize mean squared error (MSE) with the Adam optimizer on the current training window. Principals and the broker may have separate learning rates, \(\eta_{\mathrm{lr},a}\) and \(\eta_{\mathrm{lr},b}\). No explicit weight regularization is applied. At \(t = 0\), each network sets \(\mathbf W_1\) using He initialization, \(\mathbf b_1=\mathbf 0\), \(\mathbf w_2=\mathbf 0\), and \(b_2=Q\) (§2a). An untrained predictor therefore returns exactly \(Q\) for every input. Networks with seed observations then take \(E_{\mathrm{init},c}\) gradient steps on their seed histories, where \(c\in\{a,b\}\).

\textbf{Alternate ridge rule.} In the ridge experiment, principals and the broker instead use

\[
\hat q_i(\mathbf x_j)=\alpha_i+\boldsymbol\beta_i^\top\mathbf x_j,
\qquad
\hat q_b(\mathbf x_i,\mathbf x_j)=\alpha_b+\boldsymbol\gamma_b^\top\psi(\mathbf x_i,\mathbf x_j).
\]

Principal and broker fits minimize mean squared error plus \(\lambda_a\) and \(\lambda_b\), respectively, times the squared Euclidean norm of the slope coefficients. Intercepts are not penalized. Raw type and pair features are used without sample standardization. Before its first fit, a ridge predictor has zero slopes and intercept \(Q\). The reporting simulations use \(\lambda_a=\pv{settingridge_lambda_agent}\) and \(\lambda_b=\pv{settingridge_lambda_broker}\).

Subsequent NN updates warm-start from the previous weights and persistent Adam moment state. The number of gradient steps adapts to the ratio of new observations to total history:

\[E_{c,t} = \max\!\left(E_{\text{steps},c}, \; \left\lceil E_{\text{init},c} \cdot \frac{n_{\text{new}}}{n_{\text{total}}} \right\rceil\right), \qquad c\in\{a,b\}.\]

where \(n_{\text{new}}\) is the number of observations added since that model's previous fitting call, \(n_{\text{total}} = |\mathcal{H}^t|\) is full history size, and \(E_{\text{steps},c}\) is the learner-specific minimum per-call step count. For a principal, \(n_{\text{new}}\) can span multiple periods because principals fit on an alternating schedule. Principals and the broker use the same functional rule but may have different parameter values. Ridge is solved directly at each fitting call. The reporting NN values are \(\eta_{\mathrm{lr},a}=\pv{settingeta_lr_agent}\), \(\eta_{\mathrm{lr},b}=\pv{settingeta_lr_broker}\), \(E_{\mathrm{init},a}=E_{\mathrm{init},b}=\pv{settingE_init_broker}\), and \(E_{\mathrm{steps},a}=E_{\mathrm{steps},b}=\pv{settingtrain_steps_broker}\).

Initialization is period 0 for each learner. Its seed-history count is stored as the first period boundary. Each later boundary stores cumulative history at the end of a simulation period. A fitting call uses observations from the most recent \(W_p = \pv{settingtrain_window_periods}\) completed simulation periods. Before \pv{settingtrain_window_periods} simulation periods are available, the window also includes period 0. Once \pv{settingtrain_window_periods} are available, the period-0 seed history is excluded. If the selected history block exceeds \(M_{\text{train}} = \pv{settingtrain_max_obs}\) observations, fitting deterministically keeps \pv{settingtrain_max_obs} spaced observations across that block, including its oldest and newest observations. If the cap is one, it keeps the newest observation. The cap changes only the fitting sample, not the stored history or the period window. NN and ridge use the same window, cap, fitting cadence, and histories.

\subsubsection{\texorpdfstring{3b. Principal \(i\)'s Model}{3b. Principal i's Model}}\label{b.-agent-is-model}

\textbf{History.} \(\mathcal{H}_i^t = \{(\mathbf{x}_j, q_{ij})\}_{m=1}^{n_i}\) records the other party's type and realized match output from every match involving \(i\).

\textbf{Input.} The principal's network takes the partner's type as input: \(\mathbf{z} = \mathbf{x}_j\).

The hidden width is \(h_A = 2d\), so the baseline uses \pv{settingAgentWidth} hidden units.

\subsubsection{3c. Broker's Model}\label{c.-brokers-model}

\textbf{History.} \(\mathcal{H}_b^t = \{(\mathbf{x}_i, \mathbf{x}_j, q_{ij})\}_{m=1}^{n_b}\) records both parties' types and realized match output from broker-mediated matches.

\textbf{Input.} The broker's network takes a symmetric pair feature vector:

\[
\psi(\mathbf{x}_i,\mathbf{x}_j)
= \left[
\mathbf{x}_i+\mathbf{x}_j;\;
\operatorname{vech}\!\left(\frac{\mathbf{x}_i\mathbf{x}_j^\top+\mathbf{x}_j\mathbf{x}_i^\top}{2}\right)
\right].
\]

The half-vectorization \(\operatorname{vech}\) keeps the lower-triangular entries of the symmetric matrix, so the broker input dimension is \(d_B = d + d(d+1)/2\), equal to \pv{settingBrokerDimension} at \(d = \pv{settingd}\).

The broker hidden width is \(h_B = 8d\), so the baseline uses \pv{settingBrokerWidth} hidden units.

Because \(\psi(\mathbf{x}_i,\mathbf{x}_j)=\psi(\mathbf{x}_j,\mathbf{x}_i)\), each observed unordered pair contributes one training row.

The broker learns only from mediated matches. It does not observe self-search outcomes.

\subsubsection{3d. Ridge Broker Ablations}\label{d.-ridge-broker-ablations}

Three separate ridge-only ablations change the broker while leaving every principal's local ridge model unchanged. \textbf{Size-matched pooled history} uniformly samples the broker's windowed and capped fitting data without replacement, using the smaller of its available sample size and the nearest-integer median sample size among principals with nonempty histories, with half-integers rounded up. \textbf{Single-principal history} randomly retains one endpoint when each broker observation is recorded, fits \(g_b(\mathbf x)=\alpha_b+\boldsymbol\gamma_b^\top\mathbf x\), and scores a pair as \(g_b(\mathbf x_i)+g_b(\mathbf x_j)-\bar q_b\), where \(\bar q_b\) is the current fitting-sample output mean and defaults to \(Q\) before fitting. \textbf{Additive ridge} uses \(\mathbf x_i+\mathbf x_j\) without bilinear features. All use the calibrated principal and broker penalties, period window, observation cap, and fitting cadence.

\subsection{4. Match Economics}\label{match-economics}

Matches use a common reservation threshold and channel-specific cost accounting.

At initialization, each simulation computes \(\bar{q}_{\text{cal}}\) as the sample mean of \(Q+f(\mathbf{x}_i,\mathbf{x}_j)\) over \pv{settingCalibrationPairsGrouped} pairs whose indices are drawn independently and uniformly with replacement from the initial principals. This quantity scales the reservation threshold, broker fee, and self-search cost.

\subsubsection{4a. Outside Options}\label{a.-outside-options}

All principals share a reservation threshold \(r\), calibrated at initialization:

\[r = f_r \cdot \bar{q}_{\text{cal}}\]

where \(\bar{q}_{\text{cal}}\) is the calibration mean and \(f_r\) is the reservation coefficient (default \pv{settingreservation_fracFixed}). Since \(\phi\) and \(c_s\) do not depend on \(r\) (§4c), \(f_r\) may take any nonnegative value, including \(f_r \geq 1\), and is swept as an analysis axis.

\subsubsection{4b. Participation Constraints}\label{b.-participation-constraints}

Demanders emit offers only above the reservation threshold. Self-search uses the demander evaluation \(\hat{q}_{i}(\mathbf{x}_j)>r\). Brokered search uses the broker pair prediction \(\hat{q}_b(\{i,j\})>r\). A one-sided offer is accepted only if the receiver's evaluation exceeds \(r\); reciprocal offers are accepted automatically (§6d).

\subsubsection{4c. Search Costs}\label{c.-search-costs}

The self-search cost and broker fee share one calibrated level:

\[
\phi = c_s = \lambda_c \cdot \bar{q}_{\text{cal}}.
\]

\(\lambda_c\) is the shared cost rate (default \pv{settingsearch_cost_rate}). The two symbols differ in payoff accounting:

\begin{itemize}
\tightlist
\item \textbf{Self-search} costs \(c_s\) on each requested self-search position, whether or not it is filled.
\item \textbf{Brokerage} charges \(\phi\) only on successful brokered placements.
\end{itemize}

\subsection{5. Network Structure and Turnover}\label{network-structure-and-turnover}

Principals interact through a single undirected network \(G\) that determines search opportunities and structural position.

\subsubsection{5a. Network Initialization}\label{a.-network-initialization}

\(G\) is initialized as a Watts-Strogatz small-world graph on the \(N\) principals (Watts \& Strogatz, 1998), with ring degree \(k_G = \pv{settingk}\), rewiring probability \(p_{\text{rewire}} = \pv{settingp_rewire}\), and random principal order on the ring.

The broker is added as a permanent node with no type vector. Broker-node edges are synchronized to standing roster members, current-period broker clients, and principals currently engaged in broker-channel matches (§6b).

\subsubsection{5b. Match Tie Formation}\label{b.-match-tie-formation}

Each realized match adds an undirected \(i\)-\(j\) edge if one does not already exist. Ties persist until either node exits.

\subsubsection{5c. Principal Turnover}\label{c.-agent-turnover}

Principals exit independently each period with probability \(\eta\) (default \pv{settingeta}). An exiting principal is removed everywhere. Its node slot is reused by an entrant with a fresh type drawn from the same type process (§1). The entrant receives \(\lfloor k_G/2 \rfloor\) edges to existing principals sampled with probability proportional to \(\exp(-\|\mathbf{x}_{i'} - \mathbf{x}_j\|^2)\).

\subsection{6. Search}\label{search}

At the start of each period, each principal draws active relationship demand

\[d_i \sim \text{Binomial}(K,\; p_{\text{demand}}),\]

where \(K\) is the maximum number of active demands the principal can initiate in the period and \(p_{\text{demand}}\) defaults to \pv{settingp_demandFixed}. If \(d_i > 0\), the principal chooses \textbf{one channel for the batch} (§7): self-search or broker.

\subsubsection{6a. Self-Search}\label{a.-self-search}

Principal \(i\)'s self-search candidate pool has two components:

\textbf{Neighbors.} For a direct network neighbor \(j\) with prior pair-specific experience, principal \(i\) uses \(\bar{q}_{ij} = \frac{1}{n_{ij}} \sum q_{ij}^{(m)}\), where \(n_{ij}\) is the number of times \(i\) and \(j\) have matched. If the graph edge exists but \(n_{ij}=0\), the principal instead uses \(\hat q_i(\mathbf x_j)\). Initial histories seed one observation for every initial graph edge. Unobserved neighbor edges arise after turnover because entrants attach structurally before matching.

\textbf{Strangers.} Each self-searching demander \(i\) independently samples a pool \(U_i^t\) of \(\min(n_{\mathrm{strangers}}, N)\) principals uniformly without replacement, where \(n_{\mathrm{strangers}} = \pv{settingn_strangers}\) (default). The demander evaluates pool members that are not current neighbors using \(\hat{q}_i(\mathbf{x}_j)\) (§3b).

Principal \(i\) orders feasible self-search candidates by these evaluations and emits up to \(d_i\) directed offers to candidates above \(r\). Equal evaluations are ordered uniformly at random, so a principal whose predictions are constant produces a random ranking rather than an undefined ordering.

\subsubsection{6b. Broker Roster and Access}\label{b.-broker-roster-and-access}

The broker maintains a \textbf{roster} of principals it can propose as counterparties when mediating matches. In each period, it also observes the current-client set \(D^t\) of principals who outsourced. Broker access for the period is:

\[\mathcal{A}_b^t = \text{Roster}^t \cup D^t.\]

Current clients expand broker access for the period but do not become standing roster members.

\textbf{Initialization.} The roster is seeded with a target size

\[R^* = \lceil \alpha_R N \rceil,\]

where the broker roster share \(\alpha_R\) defaults to \pv{settingroster_fracFixed}. The model draws \(R^*\) principals uniformly at random (default \pv{settingRosterSize} at \(N = \pv{settingN}\)). The broker's history is seeded from existing roster-roster ties in the initial graph, capped at \pv{settingBrokerSeedCap} observed ties.

\textbf{Roster refresh.} At the start of each period, each member independently leaves the roster with probability \(p_{\text{roster}}\) (default \(\pv{settingroster_churn}\)). If \(\widetilde{\text{Roster}}^t\) is the post-churn roster,

\[\widetilde{\text{Roster}}^t = \{i \in \text{Roster}^{t-1} : u_i^t > p_{\text{roster}}\}, \qquad u_i^t \overset{iid}{\sim} \mathrm{Unif}[0,1],\]

then the broker samples without replacement from \(\{1,\ldots,N\}\setminus \widetilde{\text{Roster}}^t\) until \(|\text{Roster}^t| = R^*\), or until the population is exhausted. Broker-node edges are synchronized to the standing roster, current clients, and current broker-mediated matches (§5a).

\subsubsection{6c. Broker-Mediated Search}\label{c.-broker-mediated-search}

For every unordered pair \(\{i,j\}\) such that one side is a broker-client demander and the other side is in \(\mathcal{A}_b^t\), the broker computes predicted match value. For each broker-client demander \(i\), it ranks accessible counterparties above \(r\), breaking equal predictions uniformly at random, and keeps the highest-ranked candidates up to \(i\)'s residual active broker demand. Directed broker offers are then processed in a global score ranking with a new, independent random ordering within score ties. If both sides are broker-client demanders with remaining active demand, the same unordered pair can generate reciprocal broker offers.

\textbf{Assessment-only.} The broker scores the principal's network neighbors and an independently sampled stranger pool (§6a) using its pair predictions.

\textbf{Access-only.} Candidates come from the broker access set \(\mathcal A_b^t\). Scores use the requesting principal's pair-specific mean for experienced neighbors and its prediction model otherwise (§6a).

All other rules remain as specified. The broker still records mediated matches and updates its prediction model (§3), but its predictions do not determine access-only offers.

\subsubsection{6d. Offer Acceptance and Match Formation}\label{d.-offer-acceptance-and-match-formation}

All self-search and broker offers enter one shared market, grouped by unordered pair.

\begin{enumerate}
\def\labelenumi{\arabic{enumi}.}
\tightlist
\item If pair \(\{i,j\}\) contains reciprocal offers \(i \to j\) and \(j \to i\), the relationship is accepted automatically.
\item If pair \(\{i,j\}\) contains one directed offer \(i \to j\), receiver \(j\) evaluates \(i\) using \(\bar{q}_{ji}\) for known partners and \(\hat{q}_j(\mathbf{x}_i)\) otherwise. The offer is accepted iff this value exceeds \(r\).
\item Each accepted unordered pair forms at most one current-period relationship.
\end{enumerate}

\textbf{At market finalization, for each accepted match:} 1. Realized output is drawn: \(q_{ij} = Q + f(\mathbf{x}_i, \mathbf{x}_j) + \varepsilon_{ij}\). 2. Both parties add the observation to their histories: \(i\) adds \((\mathbf{x}_j, q_{ij})\) to \(\mathcal{H}_{i}\); \(j\) adds \((\mathbf{x}_i, q_{ij})\) to \(\mathcal{H}_{j}\). 3. If either directed offer used the broker channel, the broker adds \((\mathbf{x}_i, \mathbf{x}_j, q_{ij})\) to \(\mathcal{H}_b\) once. 4. An edge is added between \(i\) and \(j\) in \(G\) (if not already present).

A demander can have at most \(d_i\) outgoing offers, but can receive any number of acceptable incoming offers. Thus \(K\) is an active-demand bound, not a capacity constraint on receiving offers.

Before the next period begins, match lists are cleared. The next period begins with no current-period matches.

\subsection{7. The Outsourcing Decision}\label{the-outsourcing-decision}

\subsubsection{7a. Satisfaction Tracking}\label{a.-satisfaction-tracking}

Each principal \(i\) maintains a satisfaction index \(s_{i,c}^t\) for each search channel \(c \in \{\text{self}, \text{broker}\}\). For the chosen channel, the index updates once per active-demand period as an exponentially weighted moving average (EWMA) of realized match output net of search costs:

\[s_{i,c}^{t+1} = (1 - \omega)\,s_{i,c}^t + \omega \cdot \tilde{q}\]

where \(\omega = \pv{settingomega}\) and \(\tilde{q}\) is the period input in the table. The denominator is requested demand \(d_i\), so unfilled requests contribute zero output. A reciprocal relationship can update two channel-specific states, one for each directed offer.

{
\begin{longtable}[]{@{}
  >{\raggedright\arraybackslash}p{(\linewidth - 2\tabcolsep) * \real{0.2250}}
  >{\raggedright\arraybackslash}p{(\linewidth - 2\tabcolsep) * \real{0.7750}}@{}}
\toprule\noalign{}
\begin{minipage}[b]{\linewidth}\raggedright
Channel
\end{minipage} & \begin{minipage}[b]{\linewidth}\raggedright
Satisfaction input \(\tilde{q}\)
\end{minipage} \\
\midrule\noalign{}
\endhead
\bottomrule\noalign{}
\endlastfoot
Self-search & \(\dfrac{\sum q_{ij}}{d_i} - c_s\), summing over accepted directed offers made by \(i\) through self-search \\
Brokered & \(\dfrac{\sum q_{ij} - \phi \cdot n_{i,\text{broker success}}}{d_i}\), summing over accepted directed offers made by \(i\) through the broker \\
\end{longtable}
}

Brokered failure gives \(\tilde{q}=0\); self-search failure gives \(\tilde{q}=-c_s\).

\(s_{i,\text{self}}^t\) covers all self-search, including known neighbors.

At initialization, satisfaction with self-search is the mean of seed match outcomes and satisfaction with brokerage is the broker's seed-data reputation (§7c). Fresh entrants set satisfaction with self-search to the neighbors' mean for that channel and satisfaction with brokerage to current broker reputation.

The first broker choice, regardless of placement success, switches the principal from broker reputation to its own \(s_{i,\text{broker}}^t\) (§7b).

\subsubsection{7b. Decision Rule}\label{b.-decision-rule}

Each principal with \(d_i > 0\) compares:

\begin{itemize}
\tightlist
\item \textbf{\(\text{score}_{\text{self}}\)} \(= s_{i,\text{self}}^t\)
\item \textbf{\(\text{score}_{\text{broker}}\)} \(= s_{i,\text{broker}}^t\) if the principal has tried the broker, otherwise \(\text{rep}_b^t\).
\end{itemize}

The principal chooses the higher-score channel. Ties are resolved by a fair coin flip.

\subsubsection{7c. Broker Reputation}\label{c.-broker-reputation}

\[\text{rep}_b^{t+1} = \begin{cases} \frac{1}{|D^t|} \sum_{i \in D^t} s_{i,\text{broker}}^{t+1} & \text{if } D^t \neq \emptyset \\[4pt] \text{rep}_b^{t} & \text{otherwise} \end{cases}\]

where \(D^t\) is the set of principals who outsourced to the broker this period. With current clients, reputation is their mean post-update broker satisfaction; with no clients, it is unchanged. Reputation is initialized from the mean broker seed outcome.

\subsection{8. Period Measures}\label{period-measures}

Period measures are evaluated after current-period matching, satisfaction, and reputation updates, but before principal entry/exit turnover. Principal decisions do not depend on these measures.

\subsubsection{Network Measures}\label{network-measures}

Broker network-position measures use \(G\), including the broker node (§5a). They are computed when \(t \bmod \Delta_{\mathrm{net}}=0\) (default \pv{settingnetwork_measure_interval}). Other period measures are recorded each period.

\textbf{Betweenness centrality.} Freeman betweenness (Freeman, 1977) for the broker node:

\[\mathrm{BC}_b = \frac{1}{(n-1)(n-2)} \sum_{u \neq b} \sum_{v \neq u} \frac{\sigma_{uv}(b)}{\sigma_{uv}}\]

where \(n = N+1\), \(\sigma_{uv}\) is the number of shortest paths from \(u\) to \(v\), and \(\sigma_{uv}(b)\) is the number passing through \(b\). The double sum counts undirected pairs from both directions, so the denominator is \((n-1)(n-2)\) (Brandes, 2001, p.~9).

\textbf{Burt's constraint.} Broker ego-network constraint (Burt, 1992):

\[\mathrm{Constraint}_b = \sum_j \left(p_{bj} + \sum_{h \neq b,j}
p_{bh}\, p_{hj}\right)^2\]

where \(p_{bj} = 1/d_b\) and \(p_{hj} = 1/d_h\) (Everett \& Borgatti, 2020). The indirect term uses intermediary \(h\)'s degree.

\textbf{Effective size.} Non-redundant broker contacts using the Borgatti (1997) binary-undirected simplification:

\[\text{ES}_b = d_b - \frac{2t_b}{d_b},\]

where \(d_b\) is the broker's degree and \(t_b\) is the number of ties among the broker's neighbors, excluding ties to the broker.

\subsubsection{Prediction Quality}\label{prediction-quality}

Prediction quality is evaluated on random candidate pairs:

{
\begin{longtable}[]{@{}
  >{\raggedright\arraybackslash}p{(\linewidth - 4\tabcolsep) * \real{0.1951}}
  >{\raggedright\arraybackslash}p{(\linewidth - 4\tabcolsep) * \real{0.3415}}
  >{\raggedright\arraybackslash}p{(\linewidth - 4\tabcolsep) * \real{0.4634}}@{}}
\toprule\noalign{}
\begin{minipage}[b]{\linewidth}\raggedright
Sample
\end{minipage} & \begin{minipage}[b]{\linewidth}\raggedright
Construction
\end{minipage} & \begin{minipage}[b]{\linewidth}\raggedright
Reported measures
\end{minipage} \\
\midrule\noalign{}
\endhead
\bottomrule\noalign{}
\endlastfoot
Random candidate pairs & Up to \pv{settingHoldoutPrincipals} principals with nonempty histories; for each sampled principal \(i\), \(\min(\pv{settingHoldoutPartners},N-1)\) non-self partners \(j\). Each predicted match value is evaluated against \(Q + f(\mathbf{x}_i,\mathbf{x}_j)\) for the same sampled pair. Both models use the same sample, which is deterministic conditional on seed and period. & Per-principal \(R^2\), root mean squared error (RMSE), bias, and Spearman rank correlation, averaged across sampled principals. \\
\end{longtable}
}

The definitions are

\[
R^2 = 1 - \frac{\mathrm{MSE}}{\operatorname{Var}_{\mathrm{pop}}(q)}, \qquad
\mathrm{RMSE} = \sqrt{\frac{1}{n}\sum(\hat{q}-q)^2}, \qquad
\mathrm{Bias} = \frac{1}{n}\sum(\hat{q}-q).
\]

Rank correlation evaluates the ranking produced from predicted match value. Candidates are ordered by predicted match value, with equal predictions ordered uniformly at random. All tie-breaking draws are independent. For random candidate pairs, Spearman's \(\rho_S\) compares this produced ranking with ranks of underlying match value. Thus a constant predictor would produce a random ranking with expected rank correlation zero. In the definitions above, \(q\) denotes the relevant target: underlying match value for random candidate pairs. Random-candidate summaries exclude a principal's candidate set from all four measures if it contains fewer than \pv{settingMinCandidatePairsWord} pairs or if the population variance of its underlying match values is below \(\sigma_\varepsilon^2/\pv{settingVarianceDivisor}\). The \(R^2\), RMSE, bias, and rank-correlation means therefore use the same valid-principal sample.

\subsubsection{Other Measures}\label{other-measures}

{
\begin{longtable}[]{@{}
  >{\raggedright\arraybackslash}p{(\linewidth - 2\tabcolsep) * \real{0.4286}}
  >{\raggedright\arraybackslash}p{(\linewidth - 2\tabcolsep) * \real{0.5714}}@{}}
\toprule\noalign{}
\begin{minipage}[b]{\linewidth}\raggedright
Measure
\end{minipage} & \begin{minipage}[b]{\linewidth}\raggedright
Definition
\end{minipage} \\
\midrule\noalign{}
\endhead
\bottomrule\noalign{}
\endlastfoot
Access vs.~assessment & For each accepted broker-directed offer, record whether receiver \(j\) was a direct neighbor of sender \(i\) in \(G\) at the time of the offer. Non-neighbor receivers indicate access; neighbor receivers indicate assessment. \\
Realized match output by channel & Average realized match output \(\bar{q}_c^t\), where \(c \in \{\text{self}, \text{brokered}\}\). \\
Requested-position fill rates & For each channel, accepted directed offers divided by requested positions. \\
Net output per principal & Total realized output from completed relationships, counting each relationship once, minus all self-search costs and broker fees, divided by \(N\). Self-search costs apply to every self-search request, successful or not. Broker fees apply to each accepted broker-directed offer, including both offers when a reciprocal relationship fulfills two brokered requests. \\
Mean channel satisfaction & Cross-principal means \(N^{-1}\sum_i s_{i,\text{self}}^t\) and \(N^{-1}\sum_i s_{i,\text{broker}}^t\). \\
Outsourcing rate & Outsourced requested positions divided by total requested positions. \\
Broker access size & \(|\mathcal{A}_b^t| = |\text{Roster}^t \cup D^t|\), measured after outsourcing decisions form \(D^t\) and before entry/exit. \\
Counterparty concentration & Median and maximum current-period counterparties per principal. \\
Principal-node degree summaries & Mean, median, minimum, and maximum degree across principals. \\
\end{longtable}
}

\section{Part II. Baseline Values and Reproducibility}\label{part-ii.-baseline-values-and-reproducibility}

\subsection{9. Baseline Values and Reproducibility}\label{baseline-values-and-reproducibility}

\subsubsection{9a. Baseline Parameter Values}\label{a.-baseline-parameter-values}

The table lists baseline values only. Entries are arranged in two columns; each column follows the document order, with the run length listed first.

The baseline learning function is NN. Ridge and its broker variants are alternate experiments (§3).

The baseline broker provides both access and assessment. The assessment-only and access-only modes are alternate experiments (§6c).

{
\begin{center}
\small
\setlength{\tabcolsep}{3pt}
\begin{tabular}{@{}lll@{\hspace{1.6em}}lll@{}}
\toprule
Symbol & Baseline & Defined in & Symbol & Baseline & Defined in \\
\midrule
\(T\) & \pv{settingT} & §1 & \(f_r\) & \pv{settingreservation_fracFixed} & §4a \\
\(N\) & \pv{settingN} & §1 & \(\lambda_c\) & \pv{settingsearch_cost_rate} & §4c \\
\(d\) & \pv{settingd} & §1 & \(k_G\) & \pv{settingk} & §5a \\
\(d_\gamma\) & \pv{settings} & §1 & \(p_{\text{rewire}}\) & \pv{settingp_rewire} & §5a \\
\(\sigma_x\) & \pv{settingsigma_x} & §1 & \(\eta\) & \pv{settingeta} & §5c \\
\(Q\) & \pv{settingQ} & §2a & \(K\) & \pv{settingK} & §6 \\
\(\rho\) & \pv{settingrhoFixed} & §2a & \(p_{\text{demand}}\) & \pv{settingp_demandFixed} & §6 \\
\(\delta\) & \pv{settingdelta} & §2c & \(n_{\mathrm{strangers}}\) & \pv{settingn_strangers} & §6a \\
\(\sigma_\varepsilon\) & \pv{settingsigma_eps} & §2a & \(\alpha_R\) & \pv{settingroster_fracFixed} & §6b \\
\(\eta_{\mathrm{lr},a},\eta_{\mathrm{lr},b}\) & \pv{settingeta_lr_agent}, \pv{settingeta_lr_broker} & §3a & \(p_{\text{roster}}\) & \pv{settingroster_churn} & §6b \\
\(E_{\text{init},a},E_{\text{init},b}\) & \pv{settingE_init_agent}, \pv{settingE_init_broker} & §3a & \(\omega\) & \pv{settingomega} & §7a \\
\(E_{\text{steps},a},E_{\text{steps},b}\) & \pv{settingtrain_steps_agent}, \pv{settingtrain_steps_broker} & §3a & \(\Delta_{\mathrm{net}}\) & \pv{settingnetwork_measure_interval} & §8 \\
\(W_p\) & \pv{settingtrain_window_periods} periods & §3a & \(M_{\text{train}}\) & \pv{settingtrain_max_obs} & §3a \\
\bottomrule
\end{tabular}
\end{center}
}

\subsubsection{9b. Reproducibility}\label{b.-reproducibility}

All model-event randomness flows from a single integer seed. The seed determines type draws, \(G\), matching function parameters (\(\mathbf{c}\), \(\mathbf{A}\), \(\mathbf{B}\)), broker seed roster, standing-roster churn and replenishment, random ordering within market score ties, and subsequent model events. Random candidate samples and their independent principal and broker tie-breaking draws are deterministic conditional on seed and period. Simulations are fully reproducible given parameters and seed.

\section{References}\label{references}

Brandes, U. (2001). A faster algorithm for betweenness centrality. \emph{Journal of Mathematical Sociology}, \emph{25}(2), 163--177.

Borgatti, S. P. (1997). Structural holes: Unpacking Burt's redundancy measures. \emph{Connections}, \emph{20}(1), 35--38.

Burt, R. S. (1992). \emph{Structural holes: The social structure of competition}. Harvard University Press.

Everett, M. G., \& Borgatti, S. P. (2020). Unpacking Burt's constraint measure. \emph{Social Networks}, \emph{62}, 50--57.

Freeman, L. C. (1977). A set of measures of centrality based on betweenness. \emph{Sociometry}, \emph{40}(1), 35--41.

Watts, D. J., \& Strogatz, S. H. (1998). Collective dynamics of `small-world' networks. \emph{Nature}, \emph{393}(6684), 440--442.

\end{document}
